% !TEX root = ../manuscript.tex

\section*{Résumé}

Ce mémoire porte sur un projet de data science appliqué à la santé et aux modes de vie, avec pour objectif de modéliser des comportements à partir de données par apprentissage automatique. Le contexte médical s'appuie sur des seuils reconnus~: obésité (IMC $> 30$ selon l'OMS, Organisation Mondiale de la Santé~\cite{who2000}), hypercholestérolémie (cholestérol $> 200$~mg/dL selon le NCEP, National Cholesterol Education Program~\cite{ncep2002}), hypertension (pression systolique $\geq 140$~mmHg selon l'ESC, European Society of Cardiology~\cite{esc2018}). Le jeu de données utilisé est le \textit{Health \& Lifestyle Dataset}, un dataset synthétique de 100~000 individus décrits par 15 variables~: habitudes de vie (activité physique, sommeil, alimentation, tabagisme, alcool) et indicateurs biologiques (IMC, cholestérol, pression artérielle, fréquence cardiaque).

J'ai mis en place un pipeline complet selon la méthodologie CRISP-DM (Cross-Industry Standard Process for Data Mining)~\cite{chapman2000}, allant de l'exploration des données jusqu'à l'interprétabilité des modèles. L'analyse exploratoire et le feature engineering ont permis de préparer les données~: encodage des variables catégorielles, création de trois variables métier (score de style de vie composite, catégorie IMC, hypertension), standardisation. En \textbf{régression}, j'ai comparé six modèles (régression linéaire, Ridge, Lasso, ElasticNet, Random Forest, XGBoost) pour prédire le taux de cholestérol, avec tuning par \texttt{RandomizedSearchCV} et validation croisée K-Fold ($k=5$). Un data leakage a été détecté et corrigé sur la variable \texttt{hypercholesterol}~: le $R^2$ artificiel de 0.675 s'est réduit à $\approx 0$ après sa suppression. En \textbf{classification}, trois modèles ont prédit le risque de maladie (\texttt{disease\_risk}), avec deux stratégies de gestion du déséquilibre des classes comparées~: ajustement des poids (\texttt{class\_weight='balanced'}) et rééchantillonnage SMOTE~\cite{chawla2002}. En \textbf{clustering}, K-Means ($k=3$) a révélé trois profils biologiques et la CAH (liaison Ward, sous-échantillon de 5~000 individus) a isolé un groupe de fumeurs intenses hyperten dus~; les clusters ont été visualisés par ACP (24.3\% de variance en 2D). L'\textbf{interprétabilité} a été analysée via les coefficients des modèles linéaires, les feature importances et SHAP (utilisé ici comme complément visuel).

Le résultat principal est une \textbf{limite fondamentale du dataset synthétique}~: les variables ont été générées indépendamment, sans aucune relation causale. Aucun modèle n'a fait mieux qu'un prédicteur naïf ($R^2 \approx 0$, AUC $\approx 0.5$) — le Lasso annule les 16 coefficients, le Random Forest atteint $R^2 = -0.005$. La divergence entre les cinq méthodes d'interprétabilité le confirme. Ce n'est pas un échec méthodologique~: c'est la rigueur du pipeline qui a permis de le détecter. Ce travail est directement transposable à des données épidémiologiques réelles.

\bigskip
\noindent\textbf{Mots-clés~:} data science, santé, machine learning, régression, classification, clustering, interprétabilité, SHAP, SMOTE, données synthétiques, data leakage, CRISP-DM.

\newpage
